\section{Validation}
\label{sec:validation}

% =============================================================================
% 4.1  Comparison with published data
% =============================================================================
\subsection{Comparison with published resonance data}
\label{sec:published_data}

We validate the forward model against the resonance frequency
measurements reported by Yamamoto et al.~\cite{Yamamoto1980}, who
measured tap-tone frequencies for watermelons of known mass and geometry.
Additional comparison is made with the pear firmness data of De Belie et
al.~\cite{DeBelieetal2000} and the pineapple modal analysis of Chen and
De Baerdemaeker~\cite{ChenDeBaerdemaeker1993}, demonstrating that the
shell model generalises across fruit species.

% TODO: Digitise published data and compute model predictions

% TODO: Add figure
% \begin{figure}[t]
%   \centering
%   \includegraphics[width=0.7\textwidth]{figures/fig_validation}
%   \caption{Model predictions versus published experimental data.
%     (\emph{a})~Watermelon tap-tone frequencies compared with Yamamoto
%     et al.~\cite{Yamamoto1980}; the dashed line indicates perfect
%     agreement.  (\emph{b})~Pear resonance frequencies compared with
%     De Belie et al.~\cite{DeBelieetal2000}.  Error bars represent
%     reported experimental uncertainty.}
%   \label{fig:validation}
% \end{figure}

% =============================================================================
% 4.2  Sensitivity to geometric uncertainty
% =============================================================================
\subsection{Effect of geometric uncertainty}
\label{sec:geom_uncertainty}

In practice, fruit geometry is measured approximately (e.g., with
callipers or by mass--density inference).  We propagate
$\pm\SI{10}{\percent}$ uncertainty in $a$, $\zeta$, and $h$
through the forward model to quantify the resulting uncertainty in the
inverted $E_{\mathrm{rind}}$.

% TODO: Monte Carlo uncertainty propagation
